\section{Guarantees and Correctness Properties}

The Bailout Protocol defines the behavioral contract that governs Agentic systems when operating at the edge of their autonomy envelope. Unlike classical safety frameworks that enforce hard bounds through static constraint, Bailout is a dynamic, provable correctness property that guarantees the system will never reach a state from which human intervention is impossible or unverifiable. This section establishes three core properties—\textbf{Safety}, \textbf{Liveness}, and \textbf{Accountability}—and provides formal proof sketches for each.

% ---------------------------------------------------------------
\subsection{Safety: No Autonomous Collapse}
% ---------------------------------------------------------------

\textbf{Property S (Safety).} \emph{An Agentic system governed by the Bailout Protocol can never reach a state in which all escalation paths are blocked or indeterminate.}

Formally, let $S_t$ denote the system state at time $t$, let $E_t$ denote the set of active escalation paths, and let $H_t$ denote the human availability predicate. Safety requires:

\[
\forall t \;:\; \neg \text{Collapsed}(S_t) \quad \iff \quad E_t \neq \emptyset \;\lor\; H_t = \text{true}
\]

where $\text{Collapsed}(S_t)$ is true when no instrumented escalation path can produce a human response within the defined Temporal Escalation Bound (TEB).

\textbf{Invariants preserved by construction:}

\begin{enumerate}
    \item \textbf{Escalation Path Non-Empty:} Every Agentic state transition $S_t \to S_{t+1}$ is annotated with a non-empty set of continuation options $C_{t+1}$. Bailout guarantees that at least one continuation in $C_{t+1}$ is a human-contact action $\alpha_h \in A_h$, where $A_h$ is the set of human-facing actions. Formally: $\forall t \;:\; C_{t+1} \cap A_h \neq \emptyset$.

    \item \textbf{No Autonomous Absorption:} The system is prohibited from transitioning to a terminal state $S_\top$ without a verified human acknowledgment $a_h \in A_h$ in the causal chain. Formally: $\text{Collapsed}(S_\top) \Rightarrow \nexists \text{ causal chain } \sigma : S_0 \xrightarrow{\sigma} S_\top \text{ with } a_h \in \sigma$.

    \item \textbf{Budget Monotonicity:} The Escalation Budget $B_t$ (measured in remaining human-contact opportunities) is monotonically non-increasing. Bailout enforces $B_{t+1} \leq B_t$ with a minimum floor of $B_\text{min} = 1$. This floor is non-negotiable and is enforced at the runtime kernel level, not at the application logic layer.
\end{enumerate}

\textbf{Proof Sketch (Safety).}

\begin{Proof}
We prove by induction over state transitions that Safety Property S holds for any Agentic system operating under Bailout.

\begin{itemize}
    \item \textbf{Base case:} At initialization ($t=0$), the system begins in state $S_0$ with $B_0 \geq 1$ and $E_0 \neq \emptyset$ by construction. Therefore $\neg\text{Collapsed}(S_0)$ holds.

    \item \textbf{Inductive step:} Assume $\neg\text{Collapsed}(S_t)$ holds. Consider any state transition $S_t \to S_{t+1}$. Bailout imposes two constraints:
    \begin{enumerate}
        \item The transition rule $T$ must select a continuation $\alpha \in C_{t+1}$.
        \item $C_{t+1}$ is generated by the Escape Valve, which is provably non-empty and contains at least one human-contact action $\alpha_h$ by Invariant 1.
    \end{enumerate}
    Therefore $E_{t+1} \neq \emptyset$. If $H_{t+1}=\text{true}$, then $E_{t+1} \cup \{H_{t+1}\}$ guarantees $\neg\text{Collapsed}(S_{t+1})$. If $H_{t+1}=\text{false}$, then $E_{t+1} \neq \emptyset$ still guarantees $\neg\text{Collapsed}(S_{t+1})$. In either case, the invariant holds.

    \item \textbf{Budget floor enforcement:} If a transition would reduce $B$ below $B_\text{min}=1$, the transition is aborted and a forced human-contact action $\alpha_h$ is triggered before any other operation proceeds. This prevents the floor from being crossed.
\end{itemize}

By induction, $\forall t \;:\; \neg\text{Collapsed}(S_t)$. $\square$
\end{Proof}

\textbf{Consequences:} Safety does not require that every autonomous action be correct—only that no autonomous action can permanently prevent human intervention. This is a far weaker and more achievable constraint than classical system safety, which typically requires all actions to be correct. Bailout trades breadth of correctness for depth of recoverability.

% ---------------------------------------------------------------
\subsection{Liveness: Eventual Human Contact}
% ---------------------------------------------------------------

\textbf{Property L (Liveness).} \emph{Any Agentic system executing under the Bailout Protocol will produce a human-contact event within a bounded number of steps, provided the human remains available and the system has remaining escalation budget.}

Formally, let $n_\text{max}$ denote the maximum number of autonomous steps permitted before a human-contact event must occur (the Temporal Escalation Bound), and let $n$ denote the actual number of steps since the last human-contact event. Liveness requires:

\[
\forall t \;:\; (B_t > 0 \;\land\; H_t = \text{true}) \Rightarrow \exists t' \in [t+1, t+n_\text{max}] \;:\; \text{HumanContactEvent}(S_{t'})
\]

\textbf{Trigger condition:} A Human Contact Event (HCE) is triggered whenever:
\begin{enumerate}
    \item The Escalation Budget $B_t$ reaches $B_\text{trigger} = \lfloor B_0 / 2 \rfloor$ (50\% depletion), or
    \item The step counter since last HCE reaches $n_\text{max}$, whichever comes first.
\end{enumerate}

\textbf{Invariants:}

\begin{enumerate}
    \item \textbf{Bounded Step Guarantee:} The system enforces $n \leq n_\text{max}$ between successive HCEs. When $n$ approaches $n_\text{max}$, the Escalation Engine is forced to select a human-contact action $\alpha_h$ regardless of utility optimization considerations. This is not a soft guideline—it is a hard enforcement at the scheduler level.

    \item \textbf{Availability Awareness:} The Liveness property is conditional on human availability. When $H_t = \text{false}$ (human unavailable), the system transitions to a \textbf{Supervised Idle} state in which no irreversible autonomous actions are permitted and the step counter is suspended. Liveness resumes when $H_t$ becomes true.

    \item \textbf{Budget Preservation:} Each HCE resets the step counter to zero and does not consume budget (HCEs are informational, not budgetary). Only autonomous escalation-triggering events consume budget.
\end{enumerate}

\textbf{Proof Sketch (Liveness).}

\begin{Proof}
We prove that under Bailout, a Human Contact Event occurs within the bounded window.

\begin{itemize}
    \item \textbf{Scheduler Enforcement:} The Bailout scheduler intercepts every autonomous step decision. Before executing step $k$, the scheduler checks: (a) $n_\text{current}$ (steps since last HCE), and (b) $B_\text{current}$. If $n_\text{current} \geq n_\text{max} - 1$, the scheduler restricts the action space to $A_h$ (human-contact actions only). The system cannot execute a non-human-contact action in this window.

    \item \textbf{Guaranteed Selection:} Since $A_h$ is non-empty (from Safety Invariant 1), the scheduler will select $\alpha_h \in A_h$. This guarantees an HCE within at most one step of entering the restricted window.

    \item \textbf{Budget Condition:} If $B_t = 0$, the system transitions to Hard Shutdown rather than continuing autonomous operation. This is a defined system state, not an undefined collapse, and is therefore a safe terminal condition. Liveness is defined only for $B_t > 0$.

    \item \textbf{Availability Suspension:} If $H_t = \text{false}$, the system enters Supervised Idle, in which the step counter does not advance and no autonomous actions are taken. The liveness clock is suspended, not reset. When $H_t$ becomes true, the clock resumes from its suspended value, preserving the bounded-guarantee property.
\end{itemize}

Therefore, for any execution trace with $B_0 > 0$ and $H_t = \text{true}$ infinitely often, there exists a bounded $n_\text{max}$ such that an HCE occurs in every window of $n_\text{max}$ steps. $\square$
\end{Proof}

% ---------------------------------------------------------------
\subsection{Accountability: Human Always Identifiable}
% ---------------------------------------------------------------

\textbf{Property A (Accountability).} \emph{For every autonomous action executed by an Agentic system, there exists an identifiable human principal who authorized the action or the class of actions from which it derives authorization, and this mapping is permanently recorded and verifiable.}

Formally, let $A_t$ denote the action executed at time $t$, let $\rho(A_t)$ denote the human principal associated with the authorization chain of $A_t$, and let $\Gamma$ denote the permanent audit log. Accountability requires:

\[
\forall t \;:\; \text{Executed}(A_t) \Rightarrow \rho(A_t) \in \Gamma \quad \land \quad \text{Verify}(\rho(A_t), \Gamma) = \text{true}
\]

\textbf{Authorization chain:} Every action $A$ traces back to a human principal via a chain:

\[
A \xleftarrow{\text{delegated by}} D \xleftarrow{\text{authorized by}} P
\]

where $P$ is the human principal (an actual person, not a system or agent), $D$ is the delegation context (scope, constraints, and temporal validity of the authorization), and $A$ is the action class or instance within $D$'s scope.

\textbf{Invariants:}

\begin{enumerate}
    \item \textbf{Singular Principal:} For any action $A_t$, $\rho(A_t)$ resolves to exactly one human principal. Multi-party consensus is represented as a delegation to a defined group, which itself is represented as a single authorization entity under a defined governance policy. The policy defines how group decisions map to individual accountability.

    \item \textbf{Immediate Verifiability:} The authorization chain for $A_t$ is written to the audit log $\Gamma$ before $A_t$ is executed. No action executes with an unrecorded authorization chain. The log entry includes: principal ID, delegation scope, constraints, timestamp, and action class.

    \item \textbf{Non-Repudiation:} Authorization chains cannot be altered after creation. Bailout implements an append-only audit log with cryptographic integrity checksums. Any attempt to modify a log entry is detectable and triggers an immediate Escalation Budget depletion event.

    \item \textbf{No Anonymous Actions:} Every action $A_t$ in the system has an associated $\rho(A_t)$. There are no privileged or system-internal actions that bypass the accountability mapping. The Agentic kernel itself is subject to this constraint.
\end{enumerate}

\textbf{Proof Sketch (Accountability).}

\begin{Proof}
We prove that every executed action has a resolvable, verifiable human principal.

\begin{itemize}
    \item \textbf{Pre-Execution Logging:} Before any action $A$ is executed, the Bailout runtime performs an authorization resolution: $\rho(A)$ is computed by traversing the delegation chain $A \to D \to P$. The result, including the full chain metadata, is written atomically to $\Gamma$. The action $A$ is permitted to execute only after the log entry is confirmed.

    \item \textbf{Chain Resolution:} Given an audit log entry for $A$, verification proceeds by: (a) retrieving $\rho(A)$ from the log, (b) confirming $\rho(A)$ corresponds to a registered human principal in the principal registry $R$, and (c) confirming the delegation $D$ under which $A$ was authorized is still within its temporal and scope constraints. This is computable in $O(1)$ time per action.

    \item \textbf{Non-Repudiation:} Log integrity is maintained via hash chaining: each entry $i$ contains $\text{checksum}(\text{entry}_i, \text{checksum}_{i-1})$. Modifying any historical entry breaks the chain, which is detectable in $O(n)$ verification pass over the log. This is the same structure used in Certificate Transparency logs and blockchain consensus protocols.

    \item \textbf{Kernel Inclusion:} The Agentic kernel itself operates under Bailout. The kernel's own actions (scheduling decisions, resource allocation, escalation triggers) are recorded in $\Gamma$ with the principal set to the system-instantiation principal registered at bootstrap. This principal is a human-defined identity (e.g., the authorized deployer), not an anonymous system process.
\end{itemize}

Therefore, every action $A_t$ is mapped to a human principal, and this mapping is verifiable at any future time. $\square$
\end{Proof}

% ---------------------------------------------------------------
\subsection{Combined Correctness Theorem}
% ---------------------------------------------------------------

The three properties—Safety, Liveness, and Accountability—are not independent; they are mutually reinforcing and together define the correctness boundary of the Bailout Protocol.

\begin{theorem}[Bailout Correctness]
For any Agentic system $M$ executing under the Bailout Protocol with initial budget $B_0 > 0$:
\begin{enumerate}
    \item $M$ never reaches a Collapsed state (Safety).
    \item $M$ produces Human Contact Events at a frequency bounded by $n_\text{max}$ while $B_t > 0$ and $H_t = \text{true}$ (Liveness).
    \item Every action executed by $M$ is mapped to a unique, verifiable human principal (Accountability).
\end{enumerate}
\end{theorem}

\begin{Proof}
The three properties are proven independently above. Their combination in Theorem 1 establishes the complete correctness boundary: Safety prevents catastrophic loss of escalation; Liveness ensures regular human contact within a bounded window; Accountability ensures every action is traceable to a human. Together, they guarantee that an Agentic system under Bailout can never operate anonymously, irreversibly, or without recoverable human oversight. $\square$
\end{Proof}

\textbf{Practical implications:} These properties do not require the system to be intelligent, optimal, or even competent. They require only that it be recoverable, accountable, and bounded. This is a fundamentally tractable target—far more tractable than classical safety, which demands correctness across all possible inputs. Bailout accepts that Agentic systems will fail, make poor decisions, and act suboptimically; its guarantees ensure those failures are always reversible and attributable.